% Swept margin delta of the adaptive repair with four probes (Fig. sweep). The rungs are geometric, so
% that they are equally spaced on the logarithmic axis; the probe k starts at
% the rung k with the k-th radius, and the probe 2 climbs one rung because the
% box it proposed has already been solved.
\begin{tikzpicture}[x=2.1cm, y=1.0cm, font=\footnotesize]
  % Logarithmic axis of the convexity value.
  \draw[-{Stealth[length=2mm]}] (0.3,0.4) -- (0.3,4.7)
    node[above] {convexity margin $\delta$ ($\log$ scale)};
  \foreach \k in {1,...,4}{
    \draw[black!25] (0.3,\k) -- (4.6,\k);
    \node[left] at (0.3,\k) {$\delta_{\k}$};}
  \draw[black, dashed] (0.3,4) -- (4.6,4);
  \node[anchor=west, align=left] at (4.62,4)
    {$\delta_4 = \delta_{\max}$\\$= 10\times$ spread of the\\values observed};
  % Probes, each at its own rung and radius.
  \foreach \k/\radius in {1/{\Delta/2}, 2/{\Delta_2}, 3/{\Delta_3}, 4/{\Delta}}{
    \node[draw, circle, fill=blue!25, inner sep=2pt] (p\k) at (\k,\k) {};
    \node[below, align=center] at (\k,0.35)
      {probe \k\\radius $\radius$};}
  % Probe 2 proposes a solved box and climbs to the next rung.
  \node[draw, circle, fill=white, inner sep=2pt] (q2) at (2,3) {};
  \draw[-{Stealth[length=2mm]}, red!75!black, thick] (p2) -- (q2);
  \node[anchor=east, align=right, red!75!black, fill=white, inner sep=1pt] at (1.9,2.75)
    {box already solved:\\one more MILP,\\no evaluation};
  % Pairing of the radii and the rungs.
  \draw[blue!60!black, thick, densely dashed, -{Stealth[length=2mm]}]
    (1.15,0.95) -- (3.85,3.6);
  \node[anchor=north west, align=left, blue!60!black] at (2.5,1.6)
    {larger radius,\\larger margin};
\end{tikzpicture}
